\textbf{登録と判定．}開発用分割dev1で\NHsExplNCfg{}通りの構成を探索した後，構成を1つに固定し（$\lambda{=}0.95$，$g{=}1/16$，ビューの重みを次数の中央値で割った和のグラフ，クラスあたり支持12・参照2・較正2，$q{=}0.2$，折なし），ベース検出器との積の形を選ぶdev1の規則，dev2での3つの仮説，採否の規則を登録してから実行した（付録\ref{sec:app_hyp}）．dev1の規則は積を$s_0\cdot p$に決めた（$s_0\cdot s$は$\times$TINSのfar-OODが\NHsDevOneTinsFarS\%で$\sigma$の\NHsDevOneTinsSigFar\%+0.5を超える）．dev2ではH1（near-OODの\FPR{}の差が+0.5以下）が\NHsDevTwoNearD{}，H2（far-OOD）が\NHsDevTwoFarD{}，H3（$p$の誤警報率が$\alpha{=}0.05$で0.06以下，$0.2$で0.21以下）が\NHsDevTwoFprFive{}／\NHsDevTwoFprTwenty{}とすべて成立し，登録した規則では$h$を主構成にすることになった．ただし\AUC{}は$\sigma$より低く（\NHsDevTwoNearDAuc{} [\NHsDevTwoNearDAucLo, \NHsDevTwoNearDAucHi]，\NHsDevTwoFarDAuc{} [\NHsDevTwoFarDAucLo, \NHsDevTwoFarDAucHi]ポイント），\FPR{}の差の区間の上端はnear +\NRHsUpperNear{}，far +\NRHsUpperFar{}と+0.5を超えていた．規則は点推定で判定したので$h$が採用されたが，区間の上端で判定していれば採用されなかった．

\begin{table*}[htbp]
\centering
\caption{一つの量$h$（登録した構成）と交差適合した$\sigma$（単体，ベンチマークは6通りの抽出の平均，dev2は5通り；差は$h-\sigma$，抽出を単位とする対応のある95\% $t$区間）．$\times$TINSの$h$はdev1の規則で選んだ$s_0\cdot p$．}
\label{tab:h16}
\vspace{1.5mm}
{\fontsize{6.5pt}{7.6pt}\selectfont\renewcommand{\arraystretch}{1.0}%
\setlength{\tabcolsep}{1.7pt}\begin{tabular}{@{}lcc|cc|c@{}}
\hline
\rule[-2.0pt]{0pt}{8.6pt} & \multicolumn{2}{c|}{$\sigma$（交差適合）} & \multicolumn{2}{c|}{一つの量$h$} & \\
\rule[-2.0pt]{0pt}{8.6pt}設定 & \AUC & \FPR & \AUC & \FPR & \FPR{}の差 [95\%区間] \\
\hline
dev2 near（登録した確認） & 97.16 & 12.92 & 96.89 & 12.87 & -0.06 [-1.00, +0.89] \\
dev2 far（登録した確認） & 96.56 & 16.73 & 96.18 & 16.98 & +0.25 [-1.04, +1.54] \\
\hline
ImageNet-1K near-OOD & 97.94 & 11.50 & 97.18 & 13.65 & +2.14 [+1.70, +2.58] \\
ImageNet-1K far-OOD & 98.79 & 5.17 & 98.20 & 7.32 & +2.15 [+1.87, +2.42] \\
Four-OOD & 98.80 & 4.56 & 98.04 & 7.00 & +2.44 [+2.06, +2.81] \\
\hline
ImageNet-200 near-OOD & 98.17 & 9.07 & 96.57 & 8.10 & -0.97 [-1.38, -0.56] \\
ImageNet-200 far-OOD & 99.35 & 2.15 & 98.59 & 4.15 & +2.00 [+1.77, +2.23] \\
CIFAR-100 near-OOD & 99.18 & 2.90 & 98.44 & 6.42 & +3.52 [+3.17, +3.87] \\
CIFAR-100 far-OOD & 99.37 & 2.34 & 97.56 & 3.31 & +0.97 [+0.50, +1.44] \\
\hline
ImageNet-1K near-OOD $\times$TINS & 98.09 & 10.50 & 96.75 & 18.58 & +8.09 [+7.54, +8.64] \\
ImageNet-1K far-OOD $\times$TINS & 99.05 & 4.16 & 98.88 & 4.92 & +0.77 [+0.60, +0.93] \\
Four-OOD $\times$TINS & 99.09 & 3.39 & 98.97 & 3.56 & +0.18 [+0.07, +0.28] \\
\hline
\end{tabular}}%

\end{table*}

\textbf{ベンチマーク．}テスト集合では$h$は$\sigma$に及ばない（\tabref{h16}）．ImageNet-1Kの\FPR{}はnear-OODで\NHsTestNearSigFpr\%から\NHsTestNearFpr\%（差\NHsTestNearD{} [\NHsTestNearLo, \NHsTestNearHi]），far-OODで\NHsTestFarSigFpr\%から\NHsTestFarFpr\%，Four-OODで\NHsTestFourSigFpr\%から\NHsTestFourFpr\%になり，\AUC{}はそれぞれ\NRAbsHsTestNearDAucVal{}，\NRAbsHsTestFarDAucVal{}，\NRAbsHsTestFourDAucVal{}ポイント下がる．ImageNet-200のnear-OODだけは\FPR{}が\NRAbsHsInTwoHundredNearDVal{}ポイント良いが\AUC{}は\NRAbsHsInTwoHundredNearDAucVal{}ポイント低く，far-OODとCIFAR-100では悪い．$\times$TINS（$s_0\cdot p$）のnear-OODは\NHsTestNearTinsSigFpr\%から\NHsTestNearTinsFpr\%へ大きく悪化する．$p$の床の上ではベース検出器が未知の順序を決めるためで，本文がベース検出器を$\exp\sR$と掛ける理由（\secref{score}）と同じである．学習画像で較正した$p$の$\alpha{=}0.05$での誤警報率はImageNet-1Kで\NHsTestNearValFive{}／\NHsTestFarValFive{}／\NHsTestFourValFive{}であり，同じ較正の$\pfin$（\NHsTestNearSigValFive{}／\NHsTestFarSigValFive{}／\NHsTestFourSigValFive）より高い（較正画像とテスト画像は交換可能でない；\secref{fdp}）．再出現の格子では\NHsRecNCells{}条件のうち\NHsRecNBetter{}条件で$\sigma$より良いが，同じ未知クラスが1--2枚の\NHsRecNSingleton{}条件では平均\NHsRecSingletonMeanD{}ポイント悪い（NINCOの$c{=}64$，$r{=}1$で\NHsRecNincoCSixtyFourROne\%，$\sigma$は\NHsRecNincoCSixtyFourROneSig\%）．

\textbf{開発用分割から移らなかったこと．}dev2はdev1と同じ作り方で，far-OODの画像集合（OpenImage-O）も共通であり，dev2での非劣性はベンチマークに移らなかった．dev1の探索（探索時のコードの値；$\sigma$の折の割り当ては登録した実行と異なる，付録\ref{sec:app_p17}）では，支持画像と種を吸収壁にして1歩ごとの代価だけを持たせた量はnear-OODで$\sigma$に並ぶがfar-OODで大きく劣り（\FPR{} \NHsExplAbsOneNear\%／\NHsExplAbsOneFar\%，$\sigma$は\NHsExplSigNear\%／\NHsExplSigFar\%；登録した実行の$\sigma$は\NRFcDevOneNearPSixteen\%／\NRFcDevOneFarPSixteen\%），その原因は局所・大域一貫性の質量の次数成分にあった．次数だけをスコアにしても\AUC{}はfar-OODで\NHsExplDegAucFarVTwo\%（DINOv2），\NHsExplDegAucFarVThree\%（DINOv3）に達し，near-OODでは\NHsExplDegAucNearVTwo\%，\NHsExplDegAucNearVThree\%にとどまった（テスト集合では次数だけの\AUC{}はfar-OODでも\NRDecDegINFarAuroc\%で，この手がかりは移らなかった；\tabref{dec}）．時間の代価$\gamma$はこの密度の手がかりを確率の中に戻すために入れたもので，dev1では$\sigma$に並んだ（\NHsExplFinNear\%／\NHsExplFinFar\%）．ベンチマークでは再出現のない未知とfar-OODで劣ることから，$\gamma$による密度の取り込みは本文の伝播の次数成分の代わりとして十分でないと考えられる．

\textbf{系A1の実測．}再出現の格子で，$h\le\pi_Q$（\TEscape）と$\pi_Q\le\bar\phi$はすべての未知画像で成り立つ（付録\ref{sec:app_bounds}の\tabref{bounds}）．$h/\pi_Q$の中央値は\NRBdHOneHOverPiMin--\NRBdHOneHOverPiMax{}で\TEscape{}は緩くないが，\CSep{}を較正画像側の下界$\ell=\rho\omega$で読むと，$\ell$が実際のスコア$s$の\NRBdHOneEllOverSMedMin--\NRBdHOneEllOverSMedMax{}倍しかないため保証される未知は全条件で0である．実際のスコアで読めば最大\NRBdHOneCertSMax\%を保証する．下界側の弱さが，系A1が実際の検知を説明できない主な理由である．

\textbf{位置づけ．}以上から，$\sigma$を主構成のまま$h$を登録した変種として報告した（v19--v21）．これは登録からの逸脱であり，その判断はテスト集合の値を見た後に行った．ただし「結果が良ければ本文，悪ければ付録に置く」ことは実行の前に決めていた（登録文書には書いていない；付録\ref{sec:app_hyp}）．区間の上端で判定する規則なら$h$は採用されず，逸脱は生じなかった．$h$は誤警報率の保証（\TValid）と検知力（\TEscape{}--\CSep）を同じ量について述べられるが，ベンチマークでの性能は$\sigma$にも$\sR$にも及ばない．
